% =====================================================================
%  COURS DE PHYSIQUE — FICHE 4
%  Loi de la gravitation universelle
%  Document généré en LaTeX avec figures TikZ tracées « from scratch ».
%  Compilation : pdflatex (2 passes pour la table des matières).
% =====================================================================
\documentclass[11pt,a4paper]{article}

% ---------- Encodage & langue ----------
\usepackage[utf8]{inputenc}
\usepackage[T1]{fontenc}
\usepackage{lmodern}
\usepackage[french]{babel}

% ---------- Mathématiques ----------
\usepackage{amsmath,amssymb,mathtools}
\usepackage{icomma}              % virgule décimale correcte en mode maths

% ---------- Unités ----------
\usepackage{siunitx}
\sisetup{output-decimal-marker={,},per-mode=symbol}

% ---------- Couleurs, boîtes, graphismes ----------
\usepackage{xcolor}
\usepackage{colortbl}   % \rowcolor dans les tableaux
\usepackage{tikz}
\usepackage{pgfplots}
\usepackage[most]{tcolorbox}
\usepackage{enumitem}
\usepackage{array}
\usepackage{booktabs}
\usepackage{multicol}
\usepackage{fancyhdr}
\usepackage[hidelinks]{hyperref}

\usetikzlibrary{arrows.meta,positioning,calc,decorations.pathmorphing,
  decorations.markings,angles,quotes,patterns,backgrounds,
  shapes.geometric,shapes.misc,fadings,intersections,fit}
\pgfplotsset{compat=1.18}

% ---------- Géométrie de page ----------
\usepackage[a4paper,margin=2.2cm,headheight=26pt,headsep=10pt]{geometry}

% =====================================================================
%  PALETTE DE COULEURS
% =====================================================================
\definecolor{primary}{RGB}{0,151,169}      % turquoise (couleur du livre)
\definecolor{primarydark}{RGB}{0,88,101}
\definecolor{defcol}{RGB}{0,114,178}        % bleu  — définitions
\definecolor{thmcol}{RGB}{0,140,120}        % vert-bleu — théorèmes
\definecolor{formulacol}{RGB}{198,93,9}     % orange — formules
\definecolor{remarkcol}{RGB}{120,94,200}    % violet — remarques
\definecolor{attncol}{RGB}{200,40,55}       % rouge — pièges
\definecolor{excol}{RGB}{0,135,90}          % vert — exemples
\definecolor{methcol}{RGB}{90,90,100}       % gris — méthode
\definecolor{earth}{RGB}{70,110,140}        % bleu Terre
\definecolor{earthdark}{RGB}{40,70,95}
\definecolor{forcecol}{RGB}{200,40,55}      % rouge — force gravitationnelle
\definecolor{normcol}{RGB}{0,114,178}       % bleu — force normale
\definecolor{velcol}{RGB}{0,135,90}         % vert — vitesse

% =====================================================================
%  STYLE COMMUN DES BOÎTES
% =====================================================================
\tcbset{
  boxbase/.style={enhanced,breakable,boxrule=0.9pt,arc=2.5pt,
     left=3mm,right=3mm,top=2.3mm,bottom=2.3mm,
     fonttitle=\bfseries,coltitle=white,
     toptitle=0.6mm,bottomtitle=0.6mm},
}

% ---- Définition (numérotée) ----
\newtcolorbox[auto counter,number within=section]{definition}[1][]{%
  boxbase,colback=defcol!5,colframe=defcol,
  title={\textsc{Définition}~\thetcbcounter\ \ifx\relax#1\relax\relax\else\textmd{--- \itshape #1}\fi}}

% ---- Théorème / Propriété (numéroté) ----
\newtcolorbox[auto counter,number within=section]{theoreme}[1][]{%
  boxbase,colback=thmcol!6,colframe=thmcol,
  title={\textsc{Théorème / Propriété}~\thetcbcounter\ \ifx\relax#1\relax\relax\else\textmd{--- \itshape #1}\fi}}

% ---- Formule (numérotée) ----
\newtcolorbox[auto counter,number within=section]{formule}[1][]{%
  boxbase,colback=formulacol!6,colframe=formulacol,
  title={\textsc{Formule}~\thetcbcounter\ \ifx\relax#1\relax\relax\else\textmd{--- \itshape #1}\fi}}

% ---- Remarque ----
\newtcolorbox{remarque}[1][]{%
  boxbase,colback=remarkcol!6,colframe=remarkcol,
  title={\textsc{Remarque}\ifx\relax#1\relax\relax\else\textmd{ : \itshape #1}\fi}}

% ---- Attention / piège ----
\newtcolorbox{attention}[1][]{%
  boxbase,colback=attncol!6,colframe=attncol,
  title={\textsc{Attention}\ifx\relax#1\relax\relax\else\textmd{ : \itshape #1}\fi}}

% ---- Exemple ----
\newtcolorbox{exemple}[1][]{%
  boxbase,colback=excol!5,colframe=excol,
  title={\textsc{Exemple}\ifx\relax#1\relax\relax\else\textmd{ : \itshape #1}\fi}}

% ---- Méthode ----
\newtcolorbox{methode}[1][]{%
  boxbase,colback=methcol!7,colframe=methcol,sharp corners,
  title={\textsc{Méthode}\ifx\relax#1\relax\relax\else\textmd{ : \itshape #1}\fi}}

% =====================================================================
%  ENVIRONNEMENT « où : »  (légende des grandeurs d'une formule)
% =====================================================================
\newenvironment{ou}{%
  \par\smallskip\noindent\textbf{où :}\nopagebreak
  \begin{itemize}[leftmargin=1.8em,topsep=2pt,itemsep=1.5pt,parsep=0pt]}%
  {\end{itemize}}

% =====================================================================
%  EN-TÊTES ET PIEDS DE PAGE
% =====================================================================
\pagestyle{fancy}
\fancyhf{}
\fancyhead[L]{\small\textcolor{primarydark}{\textbf{Fiche 4} — Loi de la gravitation universelle}}
\fancyhead[R]{\small\textcolor{primarydark}{\textsc{Physique}}}
\fancyfoot[C]{\small\thepage}
\renewcommand{\headrulewidth}{0.8pt}
\renewcommand{\headrule}{\hbox to\headwidth{\color{primary}\leaders\hrule height \headrulewidth\hfill}}

% Titres de section en couleur
\usepackage{titlesec}
\titleformat{\section}{\Large\bfseries\color{primarydark}}{\thesection}{0.6em}{}
\titleformat{\subsection}{\large\bfseries\color{primary}}{\thesubsection}{0.6em}{}
\titleformat{\subsubsection}{\normalsize\bfseries\color{primary!80!black}}{\thesubsubsection}{0.6em}{}

% Raccourcis maths
\newcommand{\vv}[1]{\vec{#1}}
\newcommand{\dd}{\mathrm{d}}
\newcommand{\Grav}{\mathcal{G}}        % constante de gravitation
\newcommand{\Newton}{\si{\newton}}
\newcommand{\ms}{\si{\metre\per\second}}
\newcommand{\mss}{\si{\metre\per\second\squared}}

% =====================================================================
\begin{document}
\sloppy

% ---------------------------------------------------------------------
%  PAGE DE TITRE
% ---------------------------------------------------------------------
\begingroup
\thispagestyle{empty}
\centering
\vspace*{1.2cm}

\begin{tikzpicture}
\node[rectangle,rounded corners=8pt,fill=primary,text=white,
      minimum width=15.5cm,minimum height=2.6cm,align=center,inner sep=10pt]
  {{\Huge\bfseries Loi de la gravitation}\\[4pt]
   {\LARGE\bfseries universelle}};
\end{tikzpicture}

\vspace{0.4cm}
{\large\color{primarydark}\textsc{Cours de physique — Fiche n\textsuperscript{o}\,4}}

\vspace{0.7cm}

% --- Petite illustration : la Terre, un satellite et la pomme de Newton ---
\begin{tikzpicture}[scale=1.0]
  % Terre
  \shade[ball color=earth] (0,0) circle (1.4);
  \node[white] at (0,0) {\small\bfseries Terre};
  % orbite
  \draw[primary,dashed,thick] (0,0) circle (2.9);
  % satellite
  \fill[forcecol] (2.9,0) circle (0.13);
  \node[right=2pt of {(2.9,0)},forcecol,font=\scriptsize\bfseries] {satellite};
  \draw[->,>=Stealth,velcol,line width=1pt] (2.9,0)--(2.9,0.95) node[right,font=\scriptsize]{$\vv v$};
  \draw[->,>=Stealth,forcecol,line width=1pt] (2.9,0)--(2.0,0) node[midway,above,font=\scriptsize]{$\vv F$};
  % corps en chute (pomme)
  \fill[red!70!black] (-2.55,1.4) circle (0.12);
  \draw[->,>=Stealth,forcecol,line width=1pt] (-2.55,1.4)--(-1.85,0.92);
\end{tikzpicture}

\vspace{0.7cm}
\begin{minipage}{0.82\textwidth}\itshape\color{black!75}\centering
« Tout corps de l'Univers en attire un autre avec une force proportionnelle au produit
de leurs masses et inversement proportionnelle au carré de la distance qui les sépare. »\\[2pt]
\normalfont\footnotesize --- Isaac Newton, \textit{Principia Mathematica}, 1687.
\end{minipage}

\vfill
{\footnotesize\color{black!60}Cours rédigé d'après la fiche 4 du livre de physique --- figures TikZ originales.}
\vspace{1cm}
\endgroup

\newpage
% ---------------------------------------------------------------------
%  TABLE DES MATIÈRES
% ---------------------------------------------------------------------
{\color{primarydark}\hrule height 1.2pt}
\vspace{4pt}
{\Large\bfseries\color{primarydark}Table des matières}
\vspace{4pt}
{\color{primarydark}\hrule height 0.6pt}
\vspace{6pt}
\renewcommand{\contentsname}{}
\vspace{-1.2cm}
\tableofcontents
\vspace{0.4cm}
{\color{primary}\hrule height 0.6pt}

\newpage

% =====================================================================
\section{Introduction : pourquoi une « gravitation universelle » ?}
% =====================================================================

Pourquoi une pomme tombe-t-elle, alors que la Lune, elle, ne tombe jamais sur
la Terre ? Pendant des siècles, on a cru que le « ciel » et la « Terre » obéissaient
à des lois différentes. Le génie d'Isaac \textsc{Newton} (1687) fut de comprendre
que c'est \emph{la même force} qui fait tomber la pomme et qui maintient la Lune
sur son orbite : la \textbf{force de gravitation}.

Cette force est dite \emph{universelle} parce qu'elle s'exerce entre \emph{tous} les
corps possédant une masse, partout dans l'Univers : entre la Terre et vous, entre la
Terre et la Lune, entre le Soleil et les planètes, entre deux galaxies, et même
entre deux pommes posées sur une table (mais l'intensité est alors si faible qu'on
ne la ressent pas).

\begin{remarque}[ce que cette fiche va vous apprendre]
Nous allons partir de la loi de Newton, puis en déduire successivement :
\begin{itemize}[leftmargin=1.6em,topsep=2pt,itemsep=1pt]
  \item la notion de \textbf{poids} et de \textbf{pesanteur} $g$, ainsi que sa
        variation avec l'altitude ;
  \item la subtile différence entre \textbf{poids réel} et \textbf{poids apparent}
        (l'expérience de l'ascenseur) ;
  \item le mouvement des \textbf{satellites} en orbite circulaire et la notion de
        satellite \textbf{géostationnaire} ;
  \item l'\textbf{énergie} d'un satellite sur son orbite.
\end{itemize}
Chaque formule est encadrée et chacun de ses symboles est expliqué \emph{avec son
unité}. De nombreux exemples chiffrés sont entièrement détaillés.
\end{remarque}

\noindent\textbf{Notations et constantes utilisées dans tout le cours.}
Pour fixer les idées, on adopte les valeurs de référence suivantes pour la Terre
(ce sont celles utilisées dans la fiche) :

\begin{center}
\renewcommand{\arraystretch}{1.3}
\begin{tabular}{>{\bfseries}l c l}
\toprule
Grandeur & Symbole & Valeur \\
\midrule
Constante de gravitation & $\Grav$ & $\SI{6.67e-11}{\newton\metre\squared\per\kilogram\squared}$ \\
Masse de la Terre & $M_T$ & $\SI{6.0e24}{\kilogram}$ \\
Rayon de la Terre & $R_T$ & $\SI{6400}{\kilo\metre}=\SI{6.4e6}{\metre}$ \\
Pesanteur au sol & $g_0$ & $\approx \SI{9.8}{\metre\per\second\squared}$ \\
\bottomrule
\end{tabular}
\end{center}

% =====================================================================
\section{La loi de la gravitation universelle}
% =====================================================================

\subsection{Énoncé et caractéristiques des forces}

Considérons deux corps $A$ et $B$, de masses respectives $m_A$ et $m_B$, dont les
centres sont séparés d'une distance $d$. Ces deux corps exercent l'un sur l'autre
une force de gravitation.

\begin{definition}[Force de gravitation entre deux corps]
Deux corps $A$ et $B$ de masses $m_A$ et $m_B$ exercent l'un sur l'autre deux forces
de gravitation $\vv F_{A/B}$ (force de $A$ sur $B$) et $\vv F_{B/A}$ (force de $B$ sur $A$).
Ces deux forces possèdent \textbf{quatre caractéristiques} :
\begin{itemize}[leftmargin=1.6em,topsep=2pt,itemsep=1pt]
  \item elles sont \textbf{attractives} (chaque corps « tire » l'autre vers lui) ;
  \item elles sont de \textbf{sens opposés} ;
  \item elles sont \textbf{dirigées suivant la droite $(AB)$} qui relie les centres
        de gravité des deux masses ;
  \item elles ont la \textbf{même intensité} : $F_{A/B}=F_{B/A}$.
\end{itemize}
\end{definition}

% --- FIGURE 1 : deux corps A et B et leurs forces ---
\begin{center}
\begin{tikzpicture}[scale=1.0]
  % droite (AB) en pointillés
  \draw[dashed,black!55,line width=0.7pt] (-1.2,-0.8)--(9.0,1.6);
  % corps A (gauche, bas)
  \shade[ball color=black!45] (0.4,-0.5) circle (0.45);
  \fill[black] (0.4,-0.5) circle (0.05);
  \node[above left=-1pt and -2pt of {(0.4,-0.5)},font=\bfseries] {A};
  \node[below=10pt of {(0.4,-0.5)},font=\itshape] {$m_A$};
  % corps B (droite, haut)
  \shade[ball color=black!45] (7.4,1.1) circle (0.55);
  \fill[black] (7.4,1.1) circle (0.05);
  \node[above right=-2pt and -2pt of {(7.4,1.1)},font=\bfseries] {B};
  \node[below right=2pt and 0pt of {(7.4,1.1)},font=\itshape] {$m_B$};
  % force F_{B/A} : appliquée en A, dirigée vers B
  \draw[->,>=Stealth,line width=1.3pt,forcecol]
        (0.85,-0.34)--(2.7,0.3) node[midway,above,sloped,font=\small]{$\vv F_{B/A}$};
  % force F_{A/B} : appliquée en B, dirigée vers A
  \draw[->,>=Stealth,line width=1.3pt,forcecol]
        (6.9,0.92)--(5.05,0.28) node[midway,above,sloped,font=\small]{$\vv F_{A/B}$};
  % accolade "forces attractives"
  \draw[decorate,decoration={brace,amplitude=6pt},black!70]
        (1.0,0.9)--(6.7,1.95);
  \node[font=\itshape\small,black!70] at (3.7,2.0) {forces attractives};
  % distance d
  \draw[<->,>=Stealth,black!70,line width=0.7pt]
        (0.4,-1.25)--(7.4,0.35) node[midway,below,sloped,font=\itshape]{$d$};
\end{tikzpicture}\\[2pt]
{\footnotesize Figure 1 --- Les deux forces de gravitation : attractives, opposées,
portées par $(AB)$, de même intensité.}
\end{center}

\begin{attention}[« deux forces » et non « une »]
Ne confondez pas $\vv F_{A/B}$ et $\vv F_{B/A}$ : ce sont \emph{deux} forces distinctes,
qui s'appliquent sur \emph{deux} corps différents (l'une sur $B$, l'autre sur $A$).
Elles ne s'annulent donc \emph{jamais} entre elles : ce sont les deux forces d'un
couple action--réaction (3\textsuperscript{e} loi de Newton). On ne peut additionner
que des forces appliquées sur le \emph{même} corps.
\end{attention}

\subsection{L'expression mathématique de la force}

\begin{formule}[Loi de la gravitation universelle (Newton)]
L'intensité commune des deux forces vaut :
\[
\boxed{\;F_{A/B}=F_{B/A}=\Grav\times\dfrac{m_A\times m_B}{d^{\,2}}\;}
\]
\begin{ou}
  \item $F_{A/B}=F_{B/A}$ : intensité de la force de gravitation, en newtons (\si{\newton}) ;
  \item $\Grav$ : constante de gravitation universelle,
        $\Grav=\SI{6.67e-11}{\newton\metre\squared\per\kilogram\squared}$ ;
  \item $m_A,\ m_B$ : masses des deux corps, en kilogrammes (\si{\kilogram}) ;
  \item $d$ : distance entre les centres de symétrie (centres de gravité) des deux
        corps, en mètres (\si{\metre}).
\end{ou}
\end{formule}

\noindent\textbf{Comment lire cette formule ?} Elle traduit deux proportionnalités :
\begin{itemize}[leftmargin=1.6em,topsep=2pt,itemsep=1pt]
  \item $F$ est \textbf{proportionnelle au produit des masses} $m_A\times m_B$ : si on
        double l'une des masses, la force double ; si on double les deux, la force est
        multipliée par $4$.
  \item $F$ est \textbf{inversement proportionnelle au carré de la distance} $d^2$ :
        si on \emph{double} la distance, la force est divisée par $2^2=4$ ; si on
        \emph{triple} la distance, elle est divisée par $3^2=9$. C'est la fameuse
        \emph{loi en $1/d^2$}.
\end{itemize}

\begin{remarque}[unité de la constante $\Grav$]
L'unité de $\Grav$ n'est pas choisie au hasard : elle est imposée par l'équation
elle-même. En isolant $\Grav=\dfrac{F\,d^2}{m_A m_B}$, on obtient pour unité
$\dfrac{\si{\newton}\times\si{\metre\squared}}{\si{\kilogram}\times\si{\kilogram}}
 =\si{\newton\metre\squared\per\kilogram\squared}$. La valeur extrêmement petite de
$\Grav$ (de l'ordre de $10^{-11}$) explique pourquoi la gravitation est imperceptible
entre objets ordinaires et ne devient sensible que lorsqu'au moins une des masses est
\emph{astronomique} (planète, étoile\dots).
\end{remarque}

% --- Exemple détaillé 1 : deux personnes ---
\begin{exemple}[la gravitation entre deux personnes est-elle perceptible ?]
Deux amis de masses $m_A=\SI{70}{\kilogram}$ et $m_B=\SI{80}{\kilogram}$ se tiennent à
$d=\SI{1}{\metre}$ l'un de l'autre. Calculons la force de gravitation qu'ils exercent
l'un sur l'autre.

\smallskip
\textbf{Application de la formule :}
\[
F=\Grav\,\frac{m_A m_B}{d^2}
 =\SI{6.67e-11}{}\times\frac{70\times 80}{1^2}
 =\SI{6.67e-11}{}\times 5600 .
\]
\[
F \approx \SI{3.7e-7}{\newton}=\SI{0.00000037}{\newton}.
\]

\textbf{Interprétation.} Cette force vaut moins d'un \emph{millionième} de newton :
elle est totalement imperceptible (à titre de comparaison, le poids d'un cheveu est
des milliers de fois plus grand). \emph{Conclusion} : entre objets de la vie courante,
la gravitation existe bel et bien, mais elle est négligeable. Il faut une masse
gigantesque, comme celle de la Terre, pour qu'elle devienne la force dominante.
\end{exemple}

% --- Exemple détaillé 2 : Terre-Lune ---
\begin{exemple}[force d'attraction Terre--Lune]
Données : $M_T=\SI{6.0e24}{\kilogram}$, masse de la Lune
$M_L=\SI{7.3e22}{\kilogram}$, distance moyenne Terre--Lune
$d=\SI{3.84e8}{\metre}$ (soit \SI{384000}{\kilo\metre}).

\smallskip
\textbf{Étape 1 --- on écrit la formule :} $\displaystyle F=\Grav\,\frac{M_T M_L}{d^2}$.

\textbf{Étape 2 --- on calcule le numérateur :}
$\Grav\,M_T M_L=\SI{6.67e-11}{}\times \SI{6.0e24}{}\times\SI{7.3e22}{}
\approx \SI{2.92e37}{}$.

\textbf{Étape 3 --- on calcule le dénominateur :}
$d^2=(\SI{3.84e8}{})^2\approx \SI{1.47e17}{\metre\squared}$.

\textbf{Étape 4 --- on divise :}
\[
F=\frac{\SI{2.92e37}{}}{\SI{1.47e17}{}}\approx\SI{2.0e20}{\newton}.
\]

\textbf{Interprétation.} La force vaut environ $\SI{2e20}{\newton}$, soit
$200\,000\,000\,000\,000\,000\,000$ newtons ! C'est cette force colossale qui maintient
la Lune sur son orbite autour de la Terre. La même formule, les mêmes constantes :
seule la taille des masses change radicalement le résultat.
\end{exemple}

% =====================================================================
\section{Du poids à la pesanteur \texorpdfstring{$g$}{g}}
% =====================================================================

\subsection{Le poids : un cas particulier de la gravitation}

Le cas le plus important pour nous est celui où l'un des deux corps est la
\textbf{Terre}. La force que la Terre exerce sur un objet de masse $m$ situé à sa
surface (ou au-dessus), c'est ce que nous appelons couramment le \textbf{poids} de
l'objet.

% --- FIGURE 2 : la Terre et un corps à l'altitude h ---
\begin{center}
\begin{tikzpicture}[scale=1.0]
  % Terre
  \shade[ball color=earth] (0,0) circle (1.5);
  \node[white,font=\small\bfseries] at (0,-0.15) {$M_T$};
  \fill[white] (0,0) circle (0.05);
  % rayon R_T
  \draw[->,>=Stealth,white,line width=0.8pt] (0,0)--(-1.06,1.06)
        node[midway,above,sloped,font=\scriptsize,black]{$R_T$};
  % corps A à l'altitude h, sur une ligne radiale
  \coordinate (O) at (0,0);
  \coordinate (A) at (4.6,3.0);
  \fill[black] (A) circle (0.09);
  \node[above right=-1pt of A,font=\bfseries] {A};
  \node[right=1pt of A,font=\itshape] {$m$};
  % ligne radiale O->A (pointillés) = distance d
  \draw[dashed,black!60,line width=0.7pt] (O)--(A);
  % la surface coupe la ligne à distance R_T
  \coordinate (S) at ($(O)!1.5cm!(A)$);
  \fill[white] (S) circle (0.04);
  % distance d (au-dessus de la ligne)
  \node[above,sloped,font=\itshape] at ($(O)!0.55!(A)$) {$d$};
  % altitude h (portion S->A)
  \draw[<->,>=Stealth,primary,line width=0.8pt]
        ($(S)+(0.25,-0.18)$)--($(A)+(0.25,-0.18)$)
        node[midway,below,sloped,font=\itshape\small,primary]{$h$};
  % force m g appliquée en A, vers le centre
  \draw[->,>=Stealth,forcecol,line width=1.3pt]
        (A)--($(A)!1.4cm!(O)$) node[midway,above,sloped,font=\small]{$m\,\vv g$};
  % relation
  \node[primary,font=\small] at (3.4,0.3) {$d=R_T+h$};
\end{tikzpicture}\\[2pt]
{\footnotesize Figure 2 --- Un corps de masse $m$ à l'altitude $h$ : la Terre l'attire
avec la force $m\vv g$, dirigée vers le centre. La distance au centre vaut $d=R_T+h$.}
\end{center}

\begin{theoreme}[le poids est la force de gravitation terrestre]
Si l'on admet l'approximation que l'intensité de la force de gravitation exercée par la
Terre sur un corps de masse $m$ correspond au poids $P=mg$, alors :
\[
m\times g = F_{Terre/A}=\Grav\times\dfrac{M_T\times m}{d^{\,2}},
\qquad\text{avec}\qquad d=R_T+h .
\]
La masse $m$ se simplifie de part et d'autre, ce qui va nous permettre d'isoler la
pesanteur $g$ (indépendante de l'objet considéré !).
\end{theoreme}

\subsection{L'intensité de pesanteur et sa variation avec l'altitude}

En simplifiant la masse $m$ dans l'égalité $mg=\Grav\dfrac{M_T m}{d^2}$, on obtient
directement l'expression de $g$ :

\begin{formule}[Intensité de pesanteur à l'altitude $h$]
\[
\boxed{\;g_h=\dfrac{\Grav\times M_T}{(R_T+h)^2}\;}
\]
\begin{ou}
  \item $g_h$ : intensité du champ de pesanteur à l'altitude $h$, en \si{\metre\per\second\squared}
        (ou de façon équivalente en \si{\newton\per\kilogram}) ;
  \item $\Grav$ : constante de gravitation, $\SI{6.67e-11}{\newton\metre\squared\per\kilogram\squared}$ ;
  \item $M_T$ : masse de la Terre, en \si{\kilogram} ;
  \item $R_T$ : rayon de la Terre, en \si{\metre} ;
  \item $h$ : altitude du corps au-dessus de la surface, en \si{\metre}.
\end{ou}
\end{formule}

\begin{remarque}[deux unités équivalentes pour $g$]
On peut exprimer $g$ aussi bien en \si{\metre\per\second\squared} (vu comme une
\emph{accélération}, celle de la chute libre) qu'en \si{\newton\per\kilogram} (vu comme
un \emph{champ}, la force par unité de masse). Les deux sont identiques car
$\si{\newton}=\si{\kilogram}\cdot\si{\metre\per\second\squared}$, donc
$\si{\newton\per\kilogram}=\si{\metre\per\second\squared}$.
\end{remarque}

\noindent\textbf{Cas particulier : au niveau du sol $(h=0)$.} On retrouve la pesanteur
de référence $g_0$ :

\begin{formule}[Pesanteur à la surface de la Terre]
\[
\boxed{\;g_0=\dfrac{\Grav\times M_T}{R_T^{\,2}}\;}
\]
\begin{ou}
  \item $g_0$ : pesanteur au niveau du sol (altitude nulle), en \si{\metre\per\second\squared} ;
  \item les autres symboles sont identiques à la formule précédente.
\end{ou}
\end{formule}

\begin{exemple}[on retrouve $g_0\approx\SI{9.8}{\metre\per\second\squared}$ par le calcul]
Vérifions que la formule redonne bien la valeur connue de la pesanteur au sol.
\[
g_0=\frac{\Grav\,M_T}{R_T^2}
   =\frac{\SI{6.67e-11}{}\times\SI{6.0e24}{}}{(\SI{6.4e6}{})^2}.
\]
\textbf{Numérateur :} $\SI{6.67e-11}{}\times\SI{6.0e24}{}=\SI{4.00e14}{}$.\\
\textbf{Dénominateur :} $(\SI{6.4e6}{})^2=\SI{4.096e13}{}$.\\
\textbf{Quotient :}
\[
g_0=\frac{\SI{4.00e14}{}}{\SI{4.096e13}{}}\approx\SI{9.8}{\metre\per\second\squared}.
\]
\textbf{Interprétation.} On retrouve la valeur familière $g_0\approx\SI{9.8}{\metre\per\second\squared}$,
ce qui valide la cohérence du modèle. C'est cette valeur que l'on utilise pour calculer
le poids $P=mg$ d'un objet au sol.
\end{exemple}

% --- FIGURE 3 : graphe g en fonction de l'altitude ---
\begin{center}
\begin{tikzpicture}
\begin{axis}[
  width=11.5cm,height=6.4cm,
  xlabel={altitude $h$ (km)},ylabel={$g_h$ (m/s$^2$)},
  xmin=0,xmax=37000,ymin=0,ymax=10.8,
  axis lines=left,
  xtick={0,6400,12800,20000,30000,35800},
  xticklabels={$0$,$6\,400$,$12\,800$,$20\,000$,$30\,000$,$35\,800$},
  ytick={0,2,4,6,8,9.8},
  yticklabels={$0$,$2$,$4$,$6$,$8$,$9{,}8$},
  xlabel style={font=\small,at={(axis description cs:0.5,-0.16)},anchor=north},
  ylabel style={font=\small,rotate=-90,at={(axis description cs:0.07,1.04)},anchor=south},
  scaled x ticks=false,
  grid=both,grid style={black!10},
  tick label style={font=\footnotesize},
  samples=200,domain=0:37000,
]
  % g(h) = G M / (R+h)^2  avec h en km : (R+h) en m = (6400+h)*1000
  \addplot[primary,line width=1.4pt]
    {6.67e-11*6.0e24/((6400+x)*1000)^2};
  % points remarquables
  \addplot[only marks,mark=*,forcecol,mark size=1.6pt] coordinates {
    (0,9.77)(800,7.72)(1520,6.38)(35800,0.225)};
  \node[forcecol,font=\scriptsize,anchor=south west] at (axis cs:800,7.72) {sol $\to$ ISS};
  \node[forcecol,font=\scriptsize,anchor=south east] at (axis cs:35000,1.0) {orbite};
  \node[forcecol,font=\scriptsize,anchor=north east] at (axis cs:35000,1.0) {géostationnaire};
  \draw[forcecol,->,>=Stealth,line width=0.5pt] (axis cs:35000,0.85) -- (axis cs:35700,0.32);
\end{axis}
\end{tikzpicture}\\[2pt]
{\footnotesize Figure 3 --- Décroissance de la pesanteur $g$ avec l'altitude (loi en
$1/(R_T+h)^2$). Elle ne devient jamais nulle, mais diminue très vite.}
\end{center}

\begin{attention}[pourquoi $g$ ne tombe pas brutalement à zéro ?]
Beaucoup pensent qu'« en orbite, il n'y a pas de gravité ». \emph{C'est faux !} À
l'altitude de la Station spatiale internationale (\SI{400}{}–\SI{800}{\kilo\metre}),
$g$ vaut encore environ \SI{7.7}{} à \SI{8.7}{\metre\per\second\squared}, soit près de
$90\,\%$ de sa valeur au sol. Les astronautes ne flottent pas parce que la gravité
aurait disparu, mais parce qu'ils sont en \textbf{chute libre permanente} (voir
section sur les satellites). La gravité reste forte ; c'est la \emph{sensation} de poids
qui disparaît.
\end{attention}

\begin{exemple}[calculer $g$ au sommet de l'Everest puis à l'altitude d'un satellite]
\textbf{(a) Au sommet de l'Everest, $h=\SI{8.8}{\kilo\metre}=\SI{8800}{\metre}$.}
\[
g_h=\frac{\SI{4.00e14}{}}{(\SI{6.4e6}{}+\SI{8800}{})^2}
   =\frac{\SI{4.00e14}{}}{(\SI{6.4088e6}{})^2}
   =\frac{\SI{4.00e14}{}}{\SI{4.107e13}{}}\approx\SI{9.74}{\metre\per\second\squared}.
\]
La variation est minuscule (de $9{,}77$ à $9{,}74$) : à l'échelle d'une montagne,
$g$ est quasi constante. \emph{C'est pourquoi on prend $g\approx\SI{9.8}{}$ partout
au voisinage du sol.}

\smallskip
\textbf{(b) À l'altitude $h=\SI{800}{\kilo\metre}$ (orbite basse).} Ici $d=R_T+h=
\SI{6400}{}+\SI{800}{}=\SI{7200}{\kilo\metre}=\SI{7.2e6}{\metre}$, donc :
\[
g_h=\frac{\SI{4.00e14}{}}{(\SI{7.2e6}{})^2}
   =\frac{\SI{4.00e14}{}}{\SI{5.184e13}{}}\approx\SI{7.72}{\metre\per\second\squared}.
\]
\textbf{Interprétation.} Même à \SI{800}{\kilo\metre} d'altitude, $g$ ne perd qu'environ
$20\,\%$ : la gravité est encore très présente.
\end{exemple}

\subsection{Ne pas confondre masse et poids}

\begin{definition}[masse et poids]
La \textbf{masse} $m$ (en \si{\kilogram}) mesure la quantité de matière d'un corps : elle
est \emph{invariable}, identique sur la Terre, sur la Lune ou dans l'espace. \\
Le \textbf{poids} $P=mg$ (en \si{\newton}) est la \emph{force} gravitationnelle subie : il
\emph{dépend du lieu} via la valeur locale de $g$.
\end{definition}

\begin{exemple}[votre poids change de la Terre à la Lune, pas votre masse]
Un astronaute a une masse $m=\SI{80}{\kilogram}$. Sur la Lune,
$g_{\text{Lune}}\approx\SI{1.6}{\metre\per\second\squared}$ (environ $6$ fois moins que
sur Terre).
\begin{itemize}[leftmargin=1.6em,topsep=2pt,itemsep=1pt]
  \item Sur Terre : $P_T=m\,g_0=80\times 9{,}8=\SI{784}{\newton}$.
  \item Sur la Lune : $P_L=m\,g_{\text{Lune}}=80\times 1{,}6=\SI{128}{\newton}$.
\end{itemize}
\textbf{Interprétation.} Sa \emph{masse} reste $\SI{80}{\kilogram}$ partout, mais son
\emph{poids} est environ $6$ fois plus faible sur la Lune : il peut y faire des bonds
spectaculaires. C'est exactement ce que l'on observe sur les images des missions Apollo.
\end{exemple}

% =====================================================================
\section{Le poids apparent : l'expérience de l'ascenseur}
% =====================================================================

\subsection{Définition et principe}

Quand vous montez sur une balance, ce que celle-ci affiche n'est pas \emph{directement}
votre poids : c'est la force avec laquelle vous appuyez sur elle. La balance mesure en
réalité la \textbf{force normale} $N$ qu'elle exerce sur vous (qui, par la
3\textsuperscript{e} loi de Newton, est égale à celle que vous exercez sur elle).

\begin{definition}[poids apparent]
Le \textbf{poids apparent} d'un individu est la force $N$ exercée par le support sur
lequel il repose (sol, plancher, balance\dots). On ne « ressent » son poids que s'il
existe un tel support : sans support, le poids ne se ressent pas. Le poids apparent
peut être \emph{différent} du poids réel $mg$ selon l'état de mouvement.
\end{definition}

\begin{remarque}[le poids réel, lui, ne change pas]
\textbf{Point essentiel à ne jamais oublier :} la force gravitationnelle réelle $mg$
(due à l'attraction de la Terre) reste \emph{toujours la même} tant que $g$ est constante.
Ce qui change avec le mouvement, c'est le poids \emph{apparent} $N$ — la sensation de
lourdeur — et non le poids réel.
\end{remarque}

\subsection{L'analyse par la 2\texorpdfstring{$^{\text{e}}$}{e} loi de Newton}

Sur l'individu (masse $m$) dans l'ascenseur s'exercent deux forces verticales : le poids
$mg$ vers le bas et la force normale $N$ du plancher vers le haut. En orientant l'axe
vertical vers le haut, la 2\textsuperscript{e} loi de Newton s'écrit $N-mg=ma$, où $a$
est l'accélération de l'ascenseur. On en tire le poids apparent :

\begin{formule}[Poids apparent dans un ascenseur]
\[
\boxed{\;N=mg+ma=m\,(g+a)\;}
\]
\begin{ou}
  \item $N$ : poids apparent (valeur lue sur la balance), en \si{\newton} ;
  \item $m$ : masse de l'individu, en \si{\kilogram} ;
  \item $g$ : intensité de pesanteur, en \si{\metre\per\second\squared} ;
  \item $a$ : accélération algébrique de l'ascenseur (comptée $>0$ vers le haut,
        $<0$ vers le bas), en \si{\metre\per\second\squared}.
\end{ou}
\end{formule}

\noindent Cette unique formule contient \emph{tous} les cas. Il suffit de connaître le
signe de $a$ :

% --- FIGURE 4 : les quatre cas de l'ascenseur ---
\begin{center}
\begin{tikzpicture}[scale=0.95,
   cabin/.style={draw=black!55,line width=0.8pt,rounded corners=1pt},
   weight/.style={->,>=Stealth,line width=1.1pt,forcecol},
   normal/.style={->,>=Stealth,line width=1.1pt,normcol},
   person/.style={black!75,line width=1pt}]

  \foreach \x/\titre/\acc/\rel [count=\i] in {
     0/{Repos ou MRU}/{$a=0$}/{$N=mg$},
     4/{MRUA vers le haut}/{$a>0$}/{$N>mg$},
     8/{MRUA vers le bas}/{$a<0$}/{$N<mg$},
     12/{Chute libre}/{$a=-g$}/{$N=0$}}
  {
    % cabine
    \draw[cabin] (\x,0) rectangle (\x+2.6,3.0);
    % balance (rectangle au sol)
    \fill[black!15] (\x+0.7,0.05) rectangle (\x+1.9,0.35);
    \draw[black!50] (\x+0.7,0.05) rectangle (\x+1.9,0.35);
    % personne stylisée : tête + corps
    \fill[black!75] (\x+1.3,1.9) circle (0.18);
    \draw[person] (\x+1.3,1.72)--(\x+1.3,0.9);
    \draw[person] (\x+1.3,1.5)--(\x+1.0,1.15);
    \draw[person] (\x+1.3,1.5)--(\x+1.6,1.15);
    \draw[person] (\x+1.3,0.9)--(\x+1.05,0.38);
    \draw[person] (\x+1.3,0.9)--(\x+1.55,0.38);
    % poids mg (depuis le centre du corps, vers le bas) -- longueur de réf.
    \draw[weight] (\x+1.85,1.3)--(\x+1.85,0.55) node[right,font=\scriptsize,pos=0.5]{$mg$};
    % titre
    \node[font=\footnotesize\bfseries,align=center] at (\x+1.3,3.45) {\titre};
    \node[font=\scriptsize,align=center] at (\x+1.3,-0.4) {\acc};
    \node[font=\scriptsize\bfseries,forcecol,align=center] at (\x+1.3,-0.85) {\rel};
  }
  % forces normales N, ajustées par cas (depuis la balance vers le haut)
  % cas 1 : N = mg
  \draw[normal] (0.75,0.55)--(0.75,1.30) node[left,font=\scriptsize,pos=0.5]{$N$};
  % cas 2 : N > mg (plus longue)
  \draw[normal] (4.75,0.55)--(4.75,1.70) node[left,font=\scriptsize,pos=0.5]{$N$};
  \draw[->,>=Stealth,velcol,line width=0.9pt] (6.35,1.0)--(6.35,1.9) node[right,font=\scriptsize]{$\vv a$};
  % cas 3 : N < mg (plus courte)
  \draw[normal] (8.75,0.55)--(8.75,1.05) node[left,font=\scriptsize,pos=0.5]{$N$};
  \draw[->,>=Stealth,velcol,line width=0.9pt] (10.35,1.9)--(10.35,1.0) node[right,font=\scriptsize]{$\vv a$};
  % cas 4 : N = 0 (croix)
  \node[forcecol,font=\scriptsize] at (12.75,1.0) {$\times$};
  \draw[->,>=Stealth,velcol,line width=0.9pt] (14.35,1.9)--(14.35,1.0) node[right,font=\scriptsize]{$\vv a=\vv g$};
\end{tikzpicture}\\[10pt]
{\footnotesize Figure 4 --- Les quatre situations dans un ascenseur. Le poids réel $mg$
(rouge) est \emph{toujours le même} ; seul le poids apparent $N$ (bleu) varie.}
\end{center}

\begin{theoreme}[lecture des quatre cas]
\begin{itemize}[leftmargin=1.6em,topsep=2pt,itemsep=2pt]
  \item \textbf{Repos ou MRU} ($a=0$) : $N-mg=0\Rightarrow N=mg$. La balance affiche le
        poids normal.
  \item \textbf{MRUA vers le haut} ($a>0$) : $N=m(g+a)>mg$. On se sent \emph{plus lourd}
        (le plancher doit pousser plus fort pour accélérer le corps vers le haut).
  \item \textbf{MRUA vers le bas} ($a<0$) : $N=m(g+a)<mg$. On se sent \emph{plus léger}.
  \item \textbf{Chute libre} ($a=-g$) : $N=m(g-g)=0$. \textbf{Impesanteur} : la balance
        affiche $0$, on ne ressent plus aucun poids.
\end{itemize}
\end{theoreme}

\begin{exemple}[la balance dans un ascenseur qui démarre vers le haut]
Une personne de masse $m=\SI{70}{\kilogram}$ monte sur une balance dans un ascenseur.
On prend $g=\SI{9.8}{\metre\per\second\squared}$.

\smallskip
\textbf{(a) Ascenseur à l'arrêt (ou en MRU).} $a=0$ :
\[
N=m(g+a)=70\times(9{,}8+0)=\SI{686}{\newton}.
\]
La balance affiche l'équivalent de $\SI{70}{\kilogram}$ ($686/9{,}8=70$).

\smallskip
\textbf{(b) L'ascenseur démarre vers le haut} avec $a=+\SI{2}{\metre\per\second\squared}$ :
\[
N=m(g+a)=70\times(9{,}8+2)=70\times 11{,}8=\SI{826}{\newton}.
\]
La balance « voit » une masse apparente de $826/9{,}8\approx\SI{84.3}{\kilogram}$ : la
personne se sent plus lourde de $\SI{14}{\kilogram}$ ! C'est la sensation de
« tassement » qu'on éprouve au démarrage vers le haut.

\smallskip
\textbf{(c) L'ascenseur décélère / descend} avec $a=-\SI{2}{\metre\per\second\squared}$ :
\[
N=70\times(9{,}8-2)=70\times 7{,}8=\SI{546}{\newton},
\]
soit une masse apparente de $\approx\SI{55.7}{\kilogram}$. La personne se sent plus légère.

\smallskip
\textbf{(d) Le câble casse — chute libre} avec $a=-g=-\SI{9.8}{\metre\per\second\squared}$ :
\[
N=70\times(9{,}8-9{,}8)=\SI{0}{\newton}.
\]
\textbf{Interprétation.} En chute libre, la balance affiche $0$ : la personne flotte
dans la cabine. Son poids \emph{réel} vaut pourtant toujours $mg=\SI{686}{\newton}$ —
c'est exactement le mécanisme de l'impesanteur ressentie par les astronautes.
\end{exemple}

\begin{attention}[« impesanteur » ne veut pas dire « plus de gravité »]
Dans le cas (d), $g$ vaut toujours $\SI{9.8}{\metre\per\second\squared}$ et le poids réel
$mg$ est intact. Ce qui s'annule, c'est uniquement le poids \emph{apparent} $N$, parce
que la balance \emph{et} la personne tombent ensemble avec la même accélération $g$ :
il n'y a plus de force de contact entre eux.
\end{attention}

% =====================================================================
\section{Les satellites en orbite circulaire}
% =====================================================================

\subsection{Pourquoi un satellite « tombe » sans jamais toucher le sol}

Un satellite en orbite circulaire autour de la Terre est en
\textbf{mouvement circulaire uniforme} (MCU) : sa vitesse garde une valeur constante,
mais sa \emph{direction} change sans cesse. Or changer de direction, c'est accélérer :
le satellite possède une \textbf{accélération centripète}, dirigée vers le centre de la
Terre. La seule force capable de fournir cette accélération est la
\textbf{force de gravitation}. Autrement dit : \emph{la gravité joue le rôle de force
centripète}. Le satellite est en chute libre permanente — il « tombe » constamment vers
la Terre, mais sa vitesse horizontale est telle qu'il « rate » toujours le sol.

% --- FIGURE 5 : satellite en orbite circulaire ---
\begin{center}
\begin{tikzpicture}[scale=1.0]
  % Terre
  \shade[ball color=earth] (0,0) circle (1.25);
  \node[white,font=\small\bfseries] at (0,0) {$M_T$};
  \fill[white] (0,0) circle (0.04);
  % orbite
  \draw[primary,dashed,line width=1pt] (0,0) circle (2.8);
  % satellite à 40°
  \coordinate (O) at (0,0);
  \coordinate (S) at (40:2.8);
  \fill[black] (S) circle (0.10);
  \node[above right=-2pt of S,font=\bfseries\small] {satellite ($m$)};
  % rayon r = R_T + h (ligne fine noire) ; étiquette au-dessus de la ligne
  \draw[->,>=Stealth,black!70,line width=0.8pt] (O)--(S);
  \node[sloped,above,font=\itshape\small,black] at ($(O)!0.62!(S)$) {$r=R_T+h$};
  % force gravitationnelle (centripète) vers le centre ; étiquette en dessous
  \draw[->,>=Stealth,forcecol,line width=1.3pt] (S)--($(S)!1.25cm!(O)$)
        node[midway,below right=-1pt,sloped,font=\small,forcecol]{$\vv F$};
  % vitesse tangente (perpendiculaire au rayon, sens trigo)
  \draw[->,>=Stealth,velcol,line width=1.3pt] (S)--($(S)+(130:1.45)$)
        node[midway,above left=-2pt,sloped,font=\small,velcol]{$\vv v$};
\end{tikzpicture}\\[2pt]
{\footnotesize Figure 5 --- Orbite circulaire : la vitesse $\vv v$ est tangente à
l'orbite, la force de gravitation $\vv F$ (centripète) pointe vers le centre $O$ de la
Terre. Le rayon de l'orbite vaut $r=R_T+h$.}
\end{center}

\subsection{La vitesse d'un satellite}

On égale la force de gravitation et la force centripète $\frac{mv^2}{r}$ (avec
$r=R_T+h$) :
\[
\underbrace{\Grav\,\frac{M_T\,m}{(R_T+h)^2}}_{\text{force gravitationnelle}}
=\underbrace{\frac{m\,v^2}{R_T+h}}_{\text{force centripète}} .
\]
La masse $m$ du satellite se simplifie, et un facteur $(R_T+h)$ aussi. On obtient :

\begin{formule}[Vitesse d'un satellite en orbite circulaire]
\[
\boxed{\;v=\sqrt{\dfrac{\Grav\times M_T}{R_T+h}}\;}
\]
\begin{ou}
  \item $v$ : vitesse (linéaire) du satellite sur son orbite, en \si{\metre\per\second} ;
  \item $\Grav$ : constante de gravitation, $\SI{6.67e-11}{\newton\metre\squared\per\kilogram\squared}$ ;
  \item $M_T$ : masse de la Terre, en \si{\kilogram} ;
  \item $R_T+h$ : rayon de l'orbite (= rayon terrestre + altitude), en \si{\metre}.
\end{ou}
\end{formule}

\begin{remarque}[un résultat capital : $v$ ne dépend pas de la masse du satellite !]
La masse $m$ du satellite a disparu de la formule. Cela signifie qu'à une altitude
donnée, \emph{tous} les satellites — qu'ils pèsent \SI{1}{\kilogram} ou \SI{10}{\tonne} —
doivent aller à la \emph{même} vitesse pour rester sur la même orbite. Plus l'orbite est
\textbf{haute}, plus la vitesse nécessaire est \textbf{faible} ($v$ décroît quand $h$
augmente).
\end{remarque}

\begin{exemple}[vitesse et durée d'une orbite à \SI{800}{\kilo\metre} d'altitude]
On reprend $\Grav M_T=\SI{4.00e14}{}$ et $r=R_T+h=\SI{7.2e6}{\metre}$.

\smallskip
\textbf{Vitesse :}
\[
v=\sqrt{\frac{\SI{4.00e14}{}}{\SI{7.2e6}{}}}=\sqrt{\SI{5.56e7}{}}
 \approx\SI{7455}{\metre\per\second}\approx\SI{7.5}{\kilo\metre\per\second}.
\]
Un satellite à cette altitude file donc à environ \SI{27000}{\kilo\metre\per\hour} !

\smallskip
\textbf{Durée d'un tour (période) :} la circonférence de l'orbite vaut
$2\pi r=2\pi\times\SI{7.2e6}{}\approx\SI{4.52e7}{\metre}$, donc
\[
T=\frac{2\pi r}{v}=\frac{\SI{4.52e7}{}}{\SI{7455}{}}\approx\SI{6068}{\second}
 \approx\SI{101}{\minute}.
\]
\textbf{Interprétation.} Le satellite boucle un tour de la Terre en à peine plus
d'une heure et demie : c'est typiquement le cas de la Station spatiale internationale,
qui voit ainsi environ $16$ levers de Soleil par jour.
\end{exemple}

\subsection{La période de révolution — 3\texorpdfstring{$^{\text{e}}$}{e} loi de Kepler}

En remplaçant $v$ par $\dfrac{2\pi(R_T+h)}{T}$ (la vitesse = circonférence / période)
dans la relation $v^2=\dfrac{\Grav M_T}{R_T+h}$, on isole le carré de la période :

\begin{formule}[Période d'un satellite (3\textsuperscript{e} loi de Kepler)]
\[
\boxed{\;T^{2}=\dfrac{4\pi^{2}\,(R_T+h)^{3}}{\Grav\times M_T}\;}
\]
\begin{ou}
  \item $T$ : période de révolution (durée d'un tour complet), en secondes (\si{\second}) ;
  \item $R_T+h$ : rayon de l'orbite, en \si{\metre} ;
  \item $\Grav$ : constante de gravitation, en \si{\newton\metre\squared\per\kilogram\squared} ;
  \item $M_T$ : masse de la Terre, en \si{\kilogram} ;
  \item $4\pi^2$ : constante numérique sans dimension ($\approx 39{,}48$).
\end{ou}
\end{formule}

\begin{theoreme}[3\textsuperscript{e} loi de Kepler]
Le rapport $\dfrac{T^2}{(R_T+h)^3}=\dfrac{4\pi^2}{\Grav M_T}$ est une \emph{constante}
identique pour tous les satellites de la Terre. Le carré de la période est
proportionnel au cube du rayon de l'orbite : \emph{plus un satellite est loin, plus sa
période est longue} (et de façon très marquée, puisque c'est le cube qui intervient).
\end{theoreme}

\begin{remarque}[on peut aussi isoler l'altitude]
En partant de la 3\textsuperscript{e} loi de Kepler, on peut isoler l'altitude $h$ à
partir de la période $T$ imposée :
\[
h=\sqrt[3]{\dfrac{\Grav\,M_T\,T^2}{4\pi^2}}-R_T .
\]
C'est cette relation qui permet de calculer à quelle altitude placer un satellite pour
qu'il ait une période donnée — par exemple un satellite géostationnaire.
\end{remarque}

\subsection{Le satellite géostationnaire}

\begin{definition}[satellite géostationnaire]
Un satellite est \textbf{géostationnaire} s'il paraît \emph{immobile} pour un
observateur au sol. Pour cela, il faut trois conditions :
\begin{itemize}[leftmargin=1.6em,topsep=2pt,itemsep=1pt]
  \item son orbite est dans le \textbf{plan de l'équateur} ;
  \item il tourne dans le \textbf{même sens} que la Terre ;
  \item sa \textbf{période} est exactement égale à celle de rotation de la Terre sur
        elle-même ($T\approx\SI{24}{\hour}$, plus précisément \SI{23}{\hour}\,\SI{56}{\minute}).
\end{itemize}
Le satellite a alors la même \emph{vitesse angulaire} que la Terre :
$\omega_{\text{sat}}=\omega_{\text{Terre}}$.
\end{definition}

\begin{exemple}[à quelle altitude un satellite est-il géostationnaire ?]
On cherche l'altitude $h$ pour laquelle $T=\SI{86164}{\second}$ (un jour sidéral, soit
$\approx\SI{24}{\hour}$). On utilise $h=\sqrt[3]{\dfrac{\Grav M_T T^2}{4\pi^2}}-R_T$.

\smallskip
\textbf{Étape 1 --- numérateur :}
$\Grav M_T T^2=\SI{4.00e14}{}\times(\SI{86164}{})^2
=\SI{4.00e14}{}\times\SI{7.42e9}{}=\SI{2.97e24}{}$.

\textbf{Étape 2 --- on divise par $4\pi^2\approx 39{,}48$ :}
$\dfrac{\SI{2.97e24}{}}{39{,}48}=\SI{7.52e22}{}$.

\textbf{Étape 3 --- racine cubique :}
$r=\sqrt[3]{\SI{7.52e22}{}}\approx\SI{4.22e7}{\metre}=\SI{42200}{\kilo\metre}$.

\textbf{Étape 4 --- on retire le rayon terrestre :}
\[
h=r-R_T=\SI{42200}{}-\SI{6400}{}\approx\SI{35800}{\kilo\metre}.
\]
\textbf{Interprétation.} Tout satellite géostationnaire — météo, télévision, GPS relais —
se trouve nécessairement à environ \SI{35800}{\kilo\metre} d'altitude au-dessus de
l'équateur. Il n'y a qu'une seule altitude possible : c'est pourquoi cette orbite est
une « ressource » précieuse et encombrée.
\end{exemple}

\subsection{Vitesse linéaire et vitesse angulaire}

\begin{formule}[Lien entre vitesse linéaire et vitesse angulaire]
\[
\boxed{\;v=\omega\times(R_T+h)\;}
\]
\begin{ou}
  \item $v$ : vitesse linéaire du satellite, en \si{\metre\per\second} ;
  \item $\omega$ : vitesse angulaire, en \si{\radian\per\second}
        (avec $\omega=\dfrac{2\pi}{T}$) ;
  \item $R_T+h$ : rayon de l'orbite, en \si{\metre}.
\end{ou}
\end{formule}

\begin{attention}[même vitesse angulaire $\ne$ même vitesse linéaire]
Deux satellites géostationnaires (et une maison sur l'équateur !) ont la \emph{même}
vitesse angulaire $\omega$, car ils font un tour dans le même temps $T$. Mais ils n'ont
\emph{pas} la même vitesse linéaire $v$ : celle-ci dépend du rayon ($v=\omega r$). Le
satellite, plus loin du centre, va beaucoup plus vite que la maison à la surface.
\end{attention}

\begin{exemple}[classer trois satellites par vitesse]
Trois satellites sont à des altitudes $h_1>h_2>h_3$ (donc $r_1>r_2>r_3$). Comment se
classent leurs vitesses ?

\smallskip
\textbf{Vitesse linéaire} : $v=\sqrt{\dfrac{\Grav M_T}{r}}$ décroît quand $r$ augmente.
Donc $\boxed{v_3>v_2>v_1}$ : \emph{le plus bas est le plus rapide}.

\smallskip
\textbf{Vitesse angulaire} : $\omega=\dfrac{v}{r}=\sqrt{\dfrac{\Grav M_T}{r^3}}$ décroît
aussi quand $r$ augmente. Donc $\boxed{\omega_3>\omega_2>\omega_1}$ : le plus bas tourne
aussi le plus « vite » angulairement (il boucle son tour en moins de temps).

\smallskip
\textbf{Interprétation.} Pour un satellite, « plus bas » signifie à la fois plus rapide
\emph{en vitesse} et plus rapide \emph{en rotation} (période plus courte). C'est cohérent
avec la 3\textsuperscript{e} loi de Kepler.
\end{exemple}

\subsection{Comparer deux satellites : la méthode des rapports}

% --- FIGURE 6 : deux satellites à des altitudes différentes ---
\begin{center}
\begin{tikzpicture}[scale=1.05]
  % Terre
  \shade[ball color=earth] (0,0) circle (1.05);
  \node[white,font=\small\bfseries] at (0,0) {$M_T$};
  \fill[white] (0,0) circle (0.04);
  \coordinate (O) at (0,0);
  % --- satellite A (haut-gauche, plus loin) ---
  \coordinate (A) at (140:3.5);
  \draw[dashed,black!55,line width=0.7pt] (O)--(A);
  \fill[black] (A) circle (0.10);
  \node[above=1pt of A,font=\bfseries\small] {A ($m$)};
  % étiquette de distance, posée sur la ligne, fond blanc (près de la Terre)
  \node[sloped,fill=white,inner sep=1pt,font=\scriptsize] at ($(O)!0.43!(A)$)
        {$d_A=R_T+1520$ km};
  % force gravitationnelle, vers l'extérieur de la ligne
  \draw[->,>=Stealth,forcecol,line width=1.2pt] (A)--($(A)!1.15cm!(O)$)
        node[midway,above,sloped,font=\scriptsize,forcecol]{$\vv F_{T/A}$};
  % --- satellite B (droite, plus proche) ---
  \coordinate (B) at (20:3.05);
  \draw[dashed,black!55,line width=0.7pt] (O)--(B);
  \fill[black] (B) circle (0.10);
  \node[above right=0pt and -2pt of B,font=\bfseries\small] {B ($m$)};
  \node[sloped,fill=white,inner sep=1pt,font=\scriptsize] at ($(O)!0.56!(B)$)
        {$d_B=R_T+800$ km};
  \draw[->,>=Stealth,forcecol,line width=1.2pt] (B)--($(B)!1.15cm!(O)$)
        node[midway,above,sloped,font=\scriptsize,forcecol]{$\vv F_{T/B}$};
\end{tikzpicture}\\[2pt]
{\footnotesize Figure 6 --- Deux satellites de \emph{même} masse $m$ à des altitudes
différentes. Celui qui est le plus bas (B) subit la force la plus intense, car il est
plus proche du centre.}
\end{center}

Quand on compare deux situations gouvernées par la même loi en $1/d^2$, la méthode la
plus efficace est de former le \textbf{rapport} des deux forces : les constantes communes
($\Grav$, $M_T$, et même $m$ si les masses sont égales) se simplifient.

\begin{exemple}[deux satellites de même masse — comparaison complète]
Deux satellites $A$ et $B$ de masses identiques (\SI{800}{\kilogram}) sont en orbite
circulaire : $A$ à l'altitude \SI{1520}{\kilo\metre}, $B$ à \SI{800}{\kilo\metre}. On
donne $R_T=\SI{6400}{\kilo\metre}$.

\smallskip
\textbf{Étape 1 --- les distances au centre :}
\[
d_A=R_T+\SI{1520}{}=\SI{7920}{\kilo\metre},\qquad
d_B=R_T+\SI{800}{}=\SI{7200}{\kilo\metre}.
\]

\textbf{Étape 2 --- on écrit le rapport des forces.} Comme les masses sont égales,
$\Grav$, $M_T$ et $m$ se simplifient et il ne reste que les distances :
\[
\frac{F_{T/A}}{F_{T/B}}
 =\frac{\Grav\,\dfrac{M_T m}{d_A^2}}{\Grav\,\dfrac{M_T m}{d_B^2}}
 =\frac{d_B^2}{d_A^2}
 =\left(\frac{d_B}{d_A}\right)^{2}
 =\left(\frac{7200}{7920}\right)^{2}
 =\left(\frac{10}{11}\right)^{2}.
\]

\textbf{Étape 3 --- on conclut.} La racine carrée du rapport vaut
$\sqrt{F_{T/A}/F_{T/B}}=\dfrac{10}{11}\approx 0{,}91<1$. Donc $F_{T/A}<F_{T/B}$ :
\emph{le satellite le plus bas, $B$, subit bien la force la plus intense}. C'est logique :
plus on est proche du centre de la Terre, plus l'attraction est forte.

\smallskip
\textbf{Vérification par les valeurs.} Avec $\Grav M_T m=\SI{4.00e14}{}\times 800=
\SI{3.20e17}{}$ :
\[
F_{T/A}=\frac{\SI{3.20e17}{}}{(\SI{7.92e6}{})^2}\approx\SI{5103}{\newton},\qquad
F_{T/B}=\frac{\SI{3.20e17}{}}{(\SI{7.2e6}{})^2}\approx\SI{6173}{\newton}.
\]
On retrouve bien $F_{T/A}/F_{T/B}=5103/6173\approx 0{,}83=(10/11)^2$. \checkmark
\end{exemple}

\begin{exemple}[effet de « réduire l'altitude de moitié » — un piège classique]
On entend parfois : « si on réduit l'altitude de moitié, la force est multipliée par
$4$ ». \textbf{C'est faux}, et voici pourquoi. La force dépend de la distance
\emph{au centre} $d=R_T+h$, pas de la seule altitude $h$.

\smallskip
Pour que la force soit multipliée par $4$, il faudrait que $d$ (et non $h$) soit divisé
par $2$, car $F\propto 1/d^2$ et $1/(d/2)^2=4/d^2$. Or réduire l'altitude de moitié donne
$d'=R_T+\dfrac{h}{2}$, qui est \emph{bien plus grand} que $\dfrac{d}{2}=\dfrac{R_T+h}{2}$
(à cause du gros terme $R_T$). Le facteur réel d'augmentation est seulement
\[
\frac{F'}{F}=\left(\frac{R_T+h}{\,R_T+\tfrac{h}{2}\,}\right)^{2},
\]
nettement inférieur à $4$. \textbf{Interprétation.} Le rayon terrestre $R_T$ « domine »
souvent l'altitude : il ne faut jamais confondre « distance au centre » et « altitude ».
\end{exemple}

% =====================================================================
\section{L'énergie d'un satellite}
% =====================================================================

Un satellite en orbite possède deux formes d'énergie mécanique : de l'énergie
\textbf{cinétique} (car il se déplace) et de l'énergie \textbf{potentielle de gravitation}
(car il est dans le champ de la Terre).

\begin{formule}[Énergie cinétique d'un satellite]
\[
\boxed{\;E_c=\tfrac12\,m\,v^2=\dfrac{\Grav\,M_T\,m}{2\,(R_T+h)}\;}
\]
\begin{ou}
  \item $E_c$ : énergie cinétique, en joules (\si{\joule}) ;
  \item $m$ : masse du satellite, en \si{\kilogram} ;
  \item $v$ : vitesse du satellite, en \si{\metre\per\second} ;
  \item $R_T+h$ : rayon de l'orbite, en \si{\metre}
        (on a utilisé $v^2=\Grav M_T/(R_T+h)$).
\end{ou}
\end{formule}

\begin{formule}[Énergie potentielle de gravitation]
\[
\boxed{\;E_p=-\,\dfrac{\Grav\,M_T\,m}{R_T+h}\;}
\]
\begin{ou}
  \item $E_p$ : énergie potentielle de gravitation, en joules (\si{\joule}) ;
  \item le signe « $-$ » traduit le choix d'une référence nulle à l'infini : un corps
        lié à la Terre a une énergie potentielle \emph{négative} ;
  \item les autres symboles ont la même signification que ci-dessus.
\end{ou}
\end{formule}

\begin{theoreme}[énergie mécanique totale, constante sur l'orbite]
La somme $E_m=E_c+E_p$ est l'énergie mécanique. Sur une orbite circulaire :
\[
E_m=E_c+E_p=\frac{\Grav M_T m}{2(R_T+h)}-\frac{\Grav M_T m}{R_T+h}
   =-\,\frac{\Grav M_T m}{2(R_T+h)} .
\]
Cette énergie reste \textbf{constante} en tout point de l'orbite ($E_c+E_p=\text{cte}$) :
l'orbite étant circulaire, la vitesse — donc $E_c$ — et la distance — donc $E_p$ — ne
varient pas. On remarque que $E_m=-E_c$ : l'énergie totale d'un satellite lié est
négative.
\end{theoreme}

\begin{exemple}[énergies d'un satellite de \SI{800}{\kilogram} à \SI{800}{\kilo\metre}]
Reprenons $\Grav M_T=\SI{4.00e14}{}$, $m=\SI{800}{\kilogram}$, $r=\SI{7.2e6}{\metre}$,
$v\approx\SI{7455}{\metre\per\second}$.

\smallskip
\textbf{Énergie cinétique :}
\[
E_c=\tfrac12 m v^2=\tfrac12\times 800\times(7455)^2
   \approx\tfrac12\times 800\times\SI{5.56e7}{}\approx\SI{2.22e10}{\joule}.
\]
\textbf{Énergie potentielle :}
\[
E_p=-\frac{\Grav M_T m}{r}=-\frac{\SI{4.00e14}{}\times 800}{\SI{7.2e6}{}}
   \approx -\SI{4.44e10}{\joule}.
\]
\textbf{Énergie mécanique totale :}
\[
E_m=E_c+E_p\approx \SI{2.22e10}{}-\SI{4.44e10}{}=-\SI{2.22e10}{\joule}=-E_c .
\]
\textbf{Interprétation.} On vérifie bien $E_m=-E_c<0$. Le signe négatif signifie que le
satellite est \emph{lié} à la Terre : il faudrait lui fournir de l'énergie
($+\SI{2.22e10}{\joule}$) pour le « libérer » et l'envoyer à l'infini. C'est la base du
calcul des coûts énergétiques d'un lancement.
\end{exemple}

% =====================================================================
\section{Synthèse et formulaire}
% =====================================================================

\subsection{Tableau récapitulatif des formules}

\begin{center}
\renewcommand{\arraystretch}{1.7}
\begin{tabular}{>{\raggedright\arraybackslash}p{4.6cm} >{\centering\arraybackslash}p{4.6cm} >{\raggedright\arraybackslash}p{5.2cm}}
\toprule
\rowcolor{primary!15}
\textbf{Grandeur} & \textbf{Formule} & \textbf{Unités / remarque} \\
\midrule
Force de gravitation & $F=\Grav\dfrac{m_A m_B}{d^2}$ & $F$ en \si{\newton} ; $d$ = distance entre centres \\
\rowcolor{primary!7}
Poids (gravitation terrestre) & $P=mg=\Grav\dfrac{M_T m}{(R_T+h)^2}$ & $P$ en \si{\newton} \\
Pesanteur à l'altitude $h$ & $g_h=\dfrac{\Grav M_T}{(R_T+h)^2}$ & en \si{\metre\per\second\squared} ; décroît avec $h$ \\
\rowcolor{primary!7}
Pesanteur au sol & $g_0=\dfrac{\Grav M_T}{R_T^2}\approx 9{,}8$ & en \si{\metre\per\second\squared} \\
Poids apparent (ascenseur) & $N=m(g+a)$ & $a>0$ vers le haut ; chute libre $\Rightarrow N=0$ \\
\rowcolor{primary!7}
Vitesse d'un satellite & $v=\sqrt{\dfrac{\Grav M_T}{R_T+h}}$ & en \si{\metre\per\second} ; indép.\ de $m$ \\
Période (Kepler) & $T^2=\dfrac{4\pi^2(R_T+h)^3}{\Grav M_T}$ & $T$ en \si{\second} \\
\rowcolor{primary!7}
Vitesse linéaire / angulaire & $v=\omega(R_T+h)$ & $\omega=\dfrac{2\pi}{T}$ en \si{\radian\per\second} \\
Énergie mécanique & $E_m=-\dfrac{\Grav M_T m}{2(R_T+h)}=-E_c$ & en \si{\joule} ; $<0$ (corps lié) \\
\bottomrule
\end{tabular}
\end{center}

\subsection{Les dix idées-clés à retenir}

\begin{enumerate}[leftmargin=2em,topsep=3pt,itemsep=3pt]
  \item La gravitation est \textbf{universelle} : elle s'exerce entre toutes les masses,
        partout.
  \item Les deux forces de gravitation sont \emph{attractives, opposées, sur la droite
        $(AB)$, de même intensité} (couple action--réaction).
  \item Loi en $1/d^2$ : doubler la distance \emph{divise la force par $4$}.
  \item Toujours utiliser la distance \textbf{entre centres} $d=R_T+h$, jamais la seule
        altitude $h$.
  \item Le poids $P=mg$ est une \emph{force} (\si{\newton}) ; la masse $m$ (\si{\kilogram})
        est invariable.
  \item $g$ \textbf{diminue avec l'altitude} mais ne s'annule pas : en orbite basse,
        $g$ vaut encore $\sim 90\,\%$ de $g_0$.
  \item Le \textbf{poids apparent} $N=m(g+a)$ change avec le mouvement ; le poids réel
        $mg$, lui, ne change pas.
  \item En \textbf{chute libre} ($a=-g$), $N=0$ : c'est l'impesanteur, sans disparition
        de la gravité.
  \item Pour un satellite, la gravité \emph{est} la force centripète :
        $v=\sqrt{\Grav M_T/(R_T+h)}$, indépendante de sa masse ; plus haut $\Rightarrow$
        plus lent.
  \item Un satellite \textbf{géostationnaire} a $T\approx\SI{24}{\hour}$ et se trouve à
        $\approx\SI{35800}{\kilo\metre}$, dans le plan équatorial.
\end{enumerate}

\vfill
\begin{center}
\begin{tikzpicture}
\node[rectangle,rounded corners=5pt,draw=primary,line width=1pt,
      fill=primary!8,inner sep=8pt,text width=15cm,align=center]{%
  \small\color{primarydark}\textbf{En une phrase :} une \emph{seule} force — la gravitation
  en $1/d^2$ — explique à la fois la chute des corps, la valeur de $g$, la sensation de
  poids, l'orbite des satellites et leur énergie. Tout le reste n'est qu'application de
  $F=\Grav\,\dfrac{m_A m_B}{d^2}$.};
\end{tikzpicture}
\end{center}

\end{document}
